\chapter{界面详解}

MobileCAD 的界面只有四个常驻区域：\textbf{顶栏}、\textbf{特征树}、\textbf{工具条与视图区}、\textbf{元素选项卡}。
所有对话框、面板、提示都以浮层形式出现在视图区之上。理解这四块，就理解了整个应用。

\screenshot{ui-main-zh.jpg}{主界面的四个常驻区域：顶栏（1）、特征树（2）、工具条与视图区（3）、元素选项卡（4）。}

\section{顶栏}

从左到右：

\begin{itemize}
  \item \ui{☰}：主菜单（文件、编辑、视图、工具、宏、帮助）。
  \item \textbf{文档名}：当前文档（\file{.FCStd}）的名字；未保存时显示“未命名”。
  \item \textbf{副标题}：\ui{工作台 · 基于 FreeCAD 提交号}，例如“零件 · 基于 FreeCAD main e95a31c”。
        它告诉你现在处于哪个工作台、引擎对应的 FreeCAD 版本。
  \item \ui{↶ / ↷}：撤销与重做。长按（或看提示文字）会显示将要撤销的操作名，例如“撤销 凸台”。
\end{itemize}

\section{元素选项卡与工作台}

底部一排选项卡对应 FreeCAD 的工作台，但按“设计流程”重新组织过：

\begin{center}
\small
\begin{tabularx}{\linewidth}{@{}lX@{}}
\toprule
\textbf{选项卡} & \textbf{内容} \\
\midrule
\ui{零件} & 零件设计（PartDesign）+ 零件（Part）+ 草图（Sketcher）+ 曲面（Surface）+ 绘图（Draft）的组合工作区，日常建模都在这里 \\
\ui{装配} & 装配（Assembly）：关节、求解、BOM \\
\ui{工程图} & 工程图（TechDraw）：出二维图并导出 PDF/SVG/DXF \\
\ui{变量} & 电子表格（Spreadsheet）：参数表、别名、表达式 \\
\ui{更多 ▾} & BIM、FEM、CAM、Mesh、Mesh Part、Points、Reverse Engineering、Robot、OpenSCAD、Inspection、Material、Draft、Part、Part Design、Sketcher、Surface 等全部工作台 \\
\bottomrule
\end{tabularx}
\end{center}

切换工作台有两种方式：点底部选项卡；或用菜单 \menupath{视图 → 工作台} 从完整列表中选。切到“更多”里的某个
工作台后，标签会显示其名字（如“FEM ▾”）。

\section{特征树}

左侧是 \textbf{模型树}（FreeCAD 的 Tree View）。它不只是目录，而是模型的“历史与结构”：

\begin{itemize}
  \item \textbf{层级}：实体（Body）包含草图、特征；装配包含构件与关节；工程图包含页面与视图。
  \item \textbf{图标}：每种对象有对应图标（体、凸台、草图、基准面、点、孔、阵列……），与 FreeCAD 桌面一致。
  \item \textbf{状态}：红色文字表示特征出错；灰色/半透明表示隐藏；带删除线表示被压缩（Suppressed）；
        名称后带小图标表示需要重新计算（Touched）；\ui{Tip} 标记表示实体当前计算到哪一步。
      特征行右侧的小红叹号表示该特征失败，打开它的对话框会看到具体错误。
  \item \textbf{单击}：选中对象（在视图中高亮）；\textbf{双击}：打开它的对话框（实体上双击则切换“活动实体”）；
        \textbf{长按}：右键菜单。
  \item \ui{◀}：收起特征树（收起后变成 40 dp 的竖条，点一下再展开）；平板上有 14 dp 拖动条，可把树宽在
        160--480 dp 之间调整（默认 240 dp）；手机上特征树是覆盖式抽屉，从顶栏的树按钮打开，点暗色遮罩关闭。
  \item 树顶的计数（例如“15”）是当前文档的对象数。树\textbf{不支持拖拽重排}——模型结构由特征历史决定，
        想改变先后顺序请用“将 Tip 滚动到此处”和特征对话框里的引用关系。
\end{itemize}

\screenshot[0.72]{tablet-partdesign.jpg}{特征树与 PartDesign 模型。树里能看到 Pad、Groove、Mirrored、Revolution、DatumPlane、Pocket、PolarPattern 等特征的完整建模历史，这就是“参数化”的含义：每一步都可回改。}

长按对象弹出的菜单（随对象类型变化）：

\begin{center}
\small
\begin{tabularx}{\linewidth}{@{}lX@{}}
\toprule
\textbf{菜单项} & \textbf{作用} \\
\midrule
\ui{编辑 …} / \ui{编辑附着} & 打开特征对话框；或只编辑它的附着（草图/基准面挂在哪里）\\
\ui{重命名} & 修改显示名称（不改变内部类型）\\
\ui{隐藏/显示} & 切换可见性（也用于子元素）\\
\ui{压缩 / 取消压缩} & 暂时跳过该特征参与计算，用于排查错误或对比设计 \\
\ui{标记为重新计算} & 强制该对象在下次重算时更新 \\
\ui{重新计算对象} & 立即重算该对象 \\
\ui{将 Tip 滚动到此处} & 把实体的计算终点移到该特征——做“历史回放”式的设计探索 \\
\ui{移入实体} & 把外部的圆角/倒角等特征并入实体历史 \\
\ui{切换活动实体} & 多实体建模时指定新特征进入哪个实体 \\
\ui{删除} & 删除对象（可用撤销恢复）\\
\bottomrule
\end{tabularx}
\end{center}

\section{工具条与工具面板}

工具条上是一排\textbf{工具组}（toolset），例如 \ui{草图}、\ui{特征}、\ui{零件}、\ui{曲面}、\ui{Draft}。
每个工具组对应 FreeCAD 的一个或多个工具栏：

\begin{itemize}
  \item 点工具组 → 在视图区左上角展开\textbf{工具面板}：按 FreeCAD 工具栏分节排列按钮，按钮名即 FreeCAD 命令名；
  \item 面板顶部有 \ui{搜索所有工具}：输入中英文都能搜（例如“pad”或“凸台”），结果按工作台分组；
  \item 面板底部的 \ui{更多工具（N）} 会折叠/展开“该工作台不在工具栏里的命令”，保证所有命令都能触达；
  \item 工具条右侧会显示\textbf{最近使用}的 6 个工具，用竖线分隔，方便重复操作；手机上工具条可横向滚动；
  \item 再点一次工具组（或点视图空白处）关闭面板。
\end{itemize}

\screenshot{ui-toolpanel-zh.jpg}{“特征”工具面板：按 FreeCAD 工具栏分节（建模特征……），命令名与图标都来自 FreeCAD；顶部可搜索全部工具。}

\section{视图区}

\begin{itemize}
  \item \textbf{导航立方体}（右上角，112×108 dp）：六个面（上/下/前/后/右/左）对应标准视图，点面心动画转过去；
        面与面之间的边条、角块给出 26 个方向；四周的 ▲▼◀▶ 每次转 15°，↶ ↷ 转 90°；右上角的 ☰ 小菜单可切换
        \ui{正交/透视}、\ui{等轴测}、\ui{适合全部}。直接拖动立方体可自由旋转。它永远指示当前朝向。
  \item \textbf{选择}：轻触拾取面、边、顶点；已选元素在视图区底部形成一层“选择条”（标签 + ×），点 × 取消。
  \item \textbf{浮动条}：草图状态条（“完全约束 ›”）、草图数值输入、Draft 提示条、装配配合提示、选择过滤器……都
        浮在视图区边缘，不遮挡模型。
  \item \textbf{右键菜单}：视图上\textbf{双指轻点}（或鼠标右键）弹出上下文菜单，内容随当前选择变化：有面时可以
        “草图/拉伸/平面”，有边时可以“圆角/倒角”，有草图时可以“拉伸/旋转/凹槽/编辑草图”。
\end{itemize}

\section{对话框体系}

MobileCAD 有五种交互控件，认识它们就不会迷路：

\begin{description}
  \item[特征对话框（Feature Dialog）] 对应桌面 FreeCAD 的任务面板：44 dp 标题栏（▾ 折叠、✓ 确定、✗ 取消，
        两个按钮分开摆放防止误触），下方错误条，中间是属性行，底部 \ui{更多选项（N）}/\ui{收起选项}。
        属性行按 FreeCAD 各任务面板的顺序排列（例如凸台：轮廓/类型/长度/长度2/……），主行在前，其余折叠。
        行类型包括：数值（数字键盘）、枚举（≤4 项显示为分段按钮，更多则下拉）、布尔（开关）、链接
        （显示“凸台 · 边3, 边7”，点击进入拾取模式，提示 \ui{点选几何}）、颜色（色块 + 11 色调色板 + 自定义）、
        文本/列表/文件、边表（逐条边设半径）、向量/放置（每个分量一行）。\textbf{每改一项都是一次可撤销事务，
        模型实时预览}；✓ 提交，✗ 回滚到打开时的状态。
  \item[任务面板（Task Panel）] 由 FreeCAD 的 Python 命令自己构建的界面（BIM、CAM、FEM、OpenSCAD、TechDraw 等），
        被渲染成同样的触屏控件：分组框、数值框、下拉框、列表、表格、复选/单选、按钮。底部是按 FreeCAD 定义的
        \ui{确定/取消/应用/关闭}。
  \item[选择器（Chooser）] 一个标题 + 若干选项（单选或多选）。用于“选择平面”“选择关节类型”“选择要显示的视图”等。
  \item[数字键盘（Numpad）] 需要数值时从屏幕边缘弹出（平板在右侧 320 dp，手机在底部全宽）：
        键位 7 8 9 / 4 5 6 / 1 2 3 / ± 0 .，右列是退格、+、−、×、÷，底行是括号、C、✕、✓；
        单位可在顶部切换（长度 mm/cm/m/µm/in/ft，角度 °/rad）。输出就是 FreeCAD 表达式，例如 \file{25 mm}、
        \file{2 * (10 mm + 5 mm)}。属性值以 \file{=} 开头时按\textbf{表达式}处理（蓝色显示）。
  \item[文本输入] 单行或多行文本框，用于文件名、标签、表达式、宏代码等。
\end{description}

\screenshot{ui-feature-dialog-zh.jpg}{特征对话框：标题“圆锥体”，属性 角度/高度/半径 1/半径 2，底部“更多选项（6）”。左上角的箭头可折叠，右上角 ✓ 确定、✗ 取消。}

\section{偏好设置}

\menupath{编辑 → 偏好设置}（工具条上也有 \ui{偏好设置}）提供四类设置：单位系统、主题、语言、更新检查。
详见 2.6 节。

\section{平板、手机与头显}

\begin{itemize}
  \item \textbf{平板（宽 ≥ 600 dp）}：特征树常驻左侧，工具面板浮在视图左上角，对话框浮在左侧；
  \item \textbf{手机（宽 < 600 dp）}：特征树变成抽屉（点顶栏树按钮或左滑），面板从底部弹出，工具条横向滚动；
  \item \textbf{MobileCAD XR}（Quest 等头显）：整个界面渲染在一块 880×620 dp 的面板里，模型在房间中，
        数字键盘与软键盘显示为第二块浮窗；对话框改为内嵌绘制（代码中的 LocalInlineDialogs）。
\end{itemize}

\section{返回键、撤销与“卡住了”}

\begin{itemize}
  \item \textbf{返回键就是 Esc}：它按“从内到外”的顺序关闭东西——数字键盘 → 文本对话框 → 选择框 → 菜单 →
        工具面板 → Draft 工具 → 装配配合 → 任务面板 → 特征对话框 → 草图；全都关着时，应用退到后台
        \textbf{而不是退出}，避免误触丢失文档。
  \item \textbf{撤销/重做}在顶栏右侧。每个特征修改、每次属性编辑都是独立事务；撤销按钮的提示会写出具体操作名。
  \item 命令不可用时（菜单置灰、工具面板 38\% 透明），运行时点按会提示“<命令> 尚未在移动端提供”——
        这类命令通常属于桌面窗口管理，移动端没有对应物。
\end{itemize}

\section{键盘与输入法}

\begin{itemize}
  \item 需要输入文字时（表达式、文本属性、宏），软键盘自动弹出；Android 的输入法候选词可用。
  \item 连接\textbf{硬件键盘}时，输入框可直接键入；物理按键不会被输入法截获（应用为此专门处理了 Android 的 IME 焦点）。
  \item 单行输入框按回车 = 确认并收起键盘；多行输入框回车 = 换行。
\end{itemize}
